\documentclass[a4paper,12pt,leqno]{article}

\usepackage{amssymb,amsmath,amsfonts}
\usepackage[a4paper,margin=2cm]{geometry}
\usepackage{newcent}
\usepackage[scaled=1.1]{newtxmath}
\usepackage{graphicx}
\graphicspath{ {./images/} }
\usepackage[usenames,dvipsnames]{xcolor}
\usepackage{booktabs}
\usepackage{caption}
\usepackage{subcaption}

\providecommand{\tightlist}{%
  \setlength{\itemsep}{0pt}\setlength{\parskip}{0pt}}
\usepackage[pdfencoding=auto]{hyperref}

\usepackage[backend=biber, style=alphabetic]{biblatex}
\addbibresource{downscaling.bib}

% ===========================================================================
% NOTE TO SAAD (delete this block once resolved) -- three AI drafts of this
% paper were produced independently from the repo and the hackathon brief.
% They agree almost verbatim on the "Related work" section (good sign: that
% part reflects what the team actually wrote/curated) but they REPORT
% DIFFERENT NUMBERS for the same quantities (e.g. the temperature PDP range
% is variously "0.7 to 2.3", "0.8 to 2.3", "<1.0 to >2.3" across the three
% drafts). None of those numbers came from your actual notebooks -- they are
% each model's plausible-sounding guess. I have NOT picked one arbitrarily;
% every quantitative claim below is left as a [TK: ...] placeholder for you
% to fill from the real pipeline output (the report at
% https://saadsouilmi.github.io/ERA5_resid/). Please do not let a number
% survive into the final paper unless it came from that pipeline.
% ===========================================================================

\title{Where ERA5 Underestimates Temperature: A Pan-European,
Station-Level Diagnostic from GenHack'25}
\author{Saad So\underline{u}ilmi\footnote{TODO: confirm spelling/ordering of
author surname (GitHub handle is ``SaadSouilmi''; all three drafts wrote
``Soulimi'' -- please double check before submission).},
Imene Ben Rejeb-Mzah, Antoine Bezat, Anna Creti,\\
Charles-Albert Lehalle, Mengda Li, Gauthier Vermandel, Thomas Watteau}
\date{Printed the \today}

\begin{document}

\maketitle

\begin{abstract}
% TODO (Saad): tighten to ~150-200 words once the numbers below are filled
% in and the results section is final.
Rescaled reanalysis temperature fields -- gridded products such as ERA5
that assimilate, interpolate, and bias-adjust weather-station records --
are treated as ground truth by a large body of downstream research in
climate economics, energy modelling, agriculture, and impact assessment.
They are, however, only an approximation of local air temperature, and
their accuracy is neither spatially uniform nor uniform across weather
regimes. Only a handful of studies have characterised these inaccuracies,
mostly over specific regions and mostly using elevation as the sole
explanatory variable. The GenHack'25 hackathon at \'Ecole Polytechnique,
with technical support from Kayrros, gave several hundred students the
opportunity to study this discrepancy over Europe, with most teams
focusing on a specific region. In this paper we recode, harmonise, and
extend the teams' collective findings into a single, reproducible
pipeline run at the scale of the whole of Europe. Using station
temperature as the reference and the ERA5 (or rescaled) product as the
candidate field, we define station-level residuals and characterise
their structure with quantile-binned partial dependence profiles and
station-clustered bootstrap inference. We find that: (i)~on average, the
rescaled product \emph{underestimates} station temperature by roughly
[TK: X to Y]$^\circ$C; (ii)~the underestimation grows monotonically with
temperature itself; (iii)~it decreases with distance to the sea beyond a
short coastal band; (iv)~it increases with elevation, consistent with
earlier regional studies; (v)~precipitation and NDVI show at most a
minor effect at continental scale; and (vi)~wind speed strongly damps
the bias, with a pronounced temperature$\times$wind interaction whereby
calm conditions amplify the warm-temperature underestimation. We discuss
the implications of these diagnostics for statistical and
machine-learning downscaling.
\end{abstract}

\paragraph{Acknowledgment to students.} Special thanks to all GenHack'25
participating teams, whose regional analyses and empirical insights are
the raw material of this paper, and in particular to the winning teams
[TK: list team names/affiliations]. We are also grateful to Kayrros for
technical support, infrastructure, and geospatial data provision, and to
\'Ecole Polytechnique for hosting the hackathon.

% ===========================================================================
\section{Introduction}
\label{sec:intro}

A large fraction of quantitative climate research -- in climate
econometrics, energy-demand modelling, crop-yield estimation,
epidemiological exposure studies, hydrology, and extreme-event
attribution -- does not consume raw climate-model output but
\emph{rescaled} temperature fields: gridded reconstructions of
near-surface air temperature obtained by assimilating, interpolating, and
bias-adjusting weather-station records into a background field such as
ERA5 \parencite{TODO:era5_hersbach2020}. These fields are treated,
implicitly, as ground truth, either because station coverage is sparse in
the region of interest or simply because a gridded product is easier to
merge with other geospatial covariates. This substitution is convenient
but it is an approximation: ERA5 temperature is the output of a
data-assimilating physical model at $\sim$0.25\textdegree{} resolution,
not a direct observation, and there is no reason to expect its error to
be uniform in space, in time, or across weather regimes.

Surprisingly little is known about the spatial and statistical structure
of this error. A small number of studies have identified systematic
biases in mountainous terrain, where elevation and sub-grid orographic
variability degrade the match between the reanalysis grid cell and the
station it represents
\parencite{satyapragyan2026bias,sebbar2023ml_downscaling,jacobs2024filling,nogueira2022assessment,zou2022performance},
and in urban areas, where the urban heat island signal is only partially
resolved by the reanalysis grid
\parencite{jacobs2024filling,nogueira2022assessment,zou2022performance}.
These studies are regional and variable-specific: each identifies one or
two explanatory covariates (elevation, urban land cover) over one region,
and to our knowledge no study has, to date, mapped the ERA5-to-station
discrepancy jointly across the full set of plausible covariates
(temperature, precipitation, wind, distance to the sea, vegetation,
urbanisation class) at the scale of an entire continent. A systematic,
Europe-wide characterisation of \emph{where, when, and why} rescaled
temperatures depart from station measurements has been missing.

The GenHack'25 hackathon at \'Ecole Polytechnique, with technical support
from Kayrros, offered a natural route to fill part of this gap.
Several hundred students, working in teams and given paired
station--product temperature data, were asked to identify the drivers of
the discrepancy. Because the problem decomposes naturally by region, most
teams produced analyses of geographically restricted subsets, using
heterogeneous -- and sometimes mutually inconsistent -- methodologies.
Individually, these projects vary in scope, method, and rigour, as is
expected from a hackathon. Collectively, however, they converged on a
recurring, cross-team narrative: the discrepancy is not noise, it has a
sign, and a handful of covariates recur as explanatory across very
different regions. This paper takes that collective observation
seriously. We do not report on individual team submissions; instead, we
recode the teams' findings into a single, consistent diagnostic pipeline
-- partial dependence of the residual on each covariate, with
cluster-bootstrap confidence intervals resampled by station -- and re-run
it at the scale of the whole of Europe, rather than region by region, to
check which of the teams' regional findings hold up, and with what
strength, once pooled.

\paragraph{Contributions.}
\begin{enumerate}
    \item We assemble a harmonised European dataset of paired station and
    ERA5 (rescaled-product) temperatures [TK: period, number of stations,
    number of station-days], together with geospatial covariates:
    elevation, distance to the sea, NDVI, DEGURBA urbanisation class,
    precipitation, and wind speed.
    \item We define the station--product residual and characterise its
    marginal and conditional structure through quantile-binned partial
    dependence profiles, with cluster-bootstrap confidence intervals that
    respect the spatial (station-level) dependence of the data
    (Section~\ref{sec:data}).
    \item We report, at pan-European scale, the marginal association
    between the residual and six covariates -- station temperature,
    precipitation, wind speed, NDVI, distance to the sea, and elevation
    -- together with a DEGURBA urbanisation breakdown
    (Section~\ref{sec:results}): a systematic underestimation by the
    rescaled product, an amplification of the bias at high temperature
    and low wind, an elevation effect consistent with the regional
    literature, a coastal effect, and no more than a minor role for
    precipitation and vegetation at this scale.
    \item We document a temperature$\times$wind interaction whereby the
    warm-temperature underestimation is strongly amplified under low wind
    and largely damped under high wind, a result not, to our knowledge,
    previously reported at continental scale.
    \item We draw consequences for statistical and machine-learning
    downscaling, arguing that the residual structure documented here
    defines a natural target for bias-aware and residual-based
    downscaling models (Section~\ref{sec:discussion}).
\end{enumerate}

The full analysis pipeline is available as open-source code
\cite{souilmi_era5_resid_github} [TK: add Zenodo DOI if/when archived].

% ===========================================================================
\section{Related work}
\label{sec:related}

\subsection{Downscaling and rescaling: task definition}

Climate downscaling refers to the construction of local or
high-resolution climate information from coarse predictors. In the usual
supervised formulation, we observe a coarse atmospheric or land-surface
field \(x_L\) on a grid \(G_L\), optional high-resolution static or
dynamic covariates \(a_H\) on a finer grid \(G_H\), and a target field
\(y_H\) on \(G_H\). The target may be a regional climate model
simulation, a reanalysis product, a satellite product, or
station-interpolated observations. The learning problem is therefore not
merely image enlargement, but estimation of the conditional law
\[
    p(y_H \mid x_L, a_H, t, s),
\]
where \(t\) denotes time and seasonality and \(s\) denotes geographical
context. The coarse field should often be recoverable by applying a known
aggregation operator \(A\) to the high-resolution output, i.e.
\(A(y_H) \approx x_L\). This ``coarse consistency'' requirement
distinguishes climate downscaling from generic super-resolution and is
central when the output will be used for risk analysis, hydrology,
infrastructure planning, agriculture, or extreme-event attribution.

Two related but distinct tasks are commonly grouped under downscaling.
The first is \emph{statistical climate downscaling}: learning the
relationship between large-scale predictors and local climate
quantities, often with bias correction, quantile mapping, or weather
typing. The second is \emph{spatial super-resolution}: mapping a
low-resolution field to a high-resolution gridded field, frequently
conditioned on topography, land cover, vegetation indices, coast
distance, or surface radiative variables. Modern machine-learning papers
often combine both: they super-resolve a coarse field while also
correcting systematic biases and generating plausible small-scale
variability
\parencite{maraun2018statistical_downscaling,labeurthre2024overview_downscaling}.

The present paper is complementary to this literature: rather than
proposing a new downscaling method, we characterise the \emph{errors of
an existing rescaled product} that downstream studies routinely treat as
observations. In the language above, our object of study is the
discrepancy between the target \(y_H\) (station temperature) and the
rescaled estimate \(\hat y_H\) over the station network.

\subsection{Documented biases between reanalysis and station temperature}
\label{sec:related-bias}

The idea that reanalysis temperature departs from station temperature in
a structured, covariate-dependent way is not new, but the existing
evidence is fragmented across a handful of region- and variable-specific
studies. In mountainous terrain, elevation and sub-grid topographic
heterogeneity have been repeatedly flagged as drivers of the discrepancy
\parencite{satyapragyan2026bias}, and terrain-aware downscaling work over
mountain regions shows that adding elevation-related covariates
materially improves the match to station data -- itself indirect evidence
that elevation is doing real explanatory work in the raw
reanalysis-station gap \parencite{sebbar2023ml_downscaling}. In cities,
several studies report that reanalysis products under-resolve the urban
heat island and that the resulting bias correlates with built-up land
cover or urbanisation intensity
\parencite{jacobs2024filling,nogueira2022assessment,zou2022performance}.
Beyond these two covariates, we are not aware of prior work that jointly
considers wind, precipitation, vegetation, and sea proximity as
explanatory variables for the same residual, nor of prior work that does
so at continental scale with quantified, station-clustered uncertainty --
which is the gap this paper addresses.

\subsection{Modern machine-learning downscaling for climate}
\label{sec:related-ml}

\subsubsection{From statistical models to deterministic deep learning}

Classical statistical methods remain important because they make the
assumptions of the task explicit. Linear or generalized linear transfer
functions, analog methods, quantile mapping, random forests, and other
machine-learning regressors link local climate variables to large-scale
circulation and surface predictors. They are comparatively interpretable
and inexpensive, but their success depends on stationarity assumptions:
the relation learned from historical data must remain valid under future
climate states. Their treatment of spatial coherence and extremes is also
limited unless explicitly modelled. For instance, recent applications to
ERA5-Land temperature over mountains show that adding terrain-aware
covariates can substantially improve local predictions, while also
exposing the sensitivity of downscaling errors to elevation and surface
heterogeneity \parencite{sebbar2023ml_downscaling}. Broader evaluations of
statistical downscaling for climate projections similarly emphasize that
no single method dominates across variables, seasons, regions, and
extreme-event metrics
\parencite{hernanz2022statistical_downscaling,tabari2021downscaling_drought}.

Deep learning reframes downscaling as a representation-learning problem.
Early super-resolution approaches such as DeepSD adapted convolutional
neural networks to climate grids, learning an end-to-end map from coarse
fields and auxiliary topographic information to high-resolution
precipitation or temperature \parencite{vandal2017deepsd}. U-Net,
residual networks, dense blocks, and encoder-decoder architectures then
became natural choices because they mix large-scale context with local
detail. These deterministic networks are strong baselines: they are fast
at inference time, easy to train with mean-square or mean-absolute
losses, and often improve RMSE and MAE. However, regression losses
encourage conditional averages. The resulting fields may look smooth,
understate extremes, and lose fine-scale spectral power even when
pointwise scores are good. This is a key limitation for climate-risk
applications, where the distribution and spatial organization of rare
events matter as much as the conditional mean.

\subsubsection{Attention mechanisms and auxiliary geospatial information}

Attention mechanisms address a specific weakness of plain convolutional
downscalers: not all covariates are equally informative everywhere. In
land surface temperature downscaling, for example, vegetation, water,
built-up area, slope, aspect, and surface reflectance may explain
different parts of the temperature signal depending on season and local
geography. The Attention Mechanism U-Net (AMUN) of
\textcite{li2025attention_downscaling} is a useful example. The method
uses Google Earth Engine to assemble ERA5-Land LST and seven
high-resolution auxiliary factors over China, then learns to downscale
hourly monthly-mean LST from \(0.1^\circ\) to \(0.01^\circ\). Its Global
Multi-Factor Cross-Attention module recalibrates the interaction between
coarse LST and auxiliary features; a feature-fusion residual dense block
extracts multilevel LST features; and a U-Net decoder reconstructs the
fine field. In simulated and real experiments, the model reports
substantially lower RMSE than a random-forest baseline.

This line of work is relevant beyond LST because it points to a general
design principle: downscaling models should learn \emph{where} and
\emph{when} to trust each explanatory variable. Orography may dominate
mountain temperature gradients, NDVI and land cover may dominate surface
heat contrasts, and large-scale humidity and winds may control
precipitation organization. Attention layers, cross-attention between
low-resolution climate states and high-resolution covariate maps, and
channel-wise gating are therefore natural tools for geophysical
downscaling. Their limitation is that they typically produce a single
best estimate unless combined with a probabilistic output layer or a
generative model. Our own diagnostic (Section~\ref{sec:results}) can be
read as the descriptive counterpart to this design principle: before
choosing which covariates a downscaling network should attend to, one
should establish which covariates actually correlate with the raw
reanalysis-station gap, and by how much -- which is exactly what the
partial-dependence analysis below provides.

\subsubsection{Generative downscaling: GANs, diffusion, and consistency
models}

Generative models are attractive for climate downscaling because the
target is intrinsically one-to-many. A coarse grid cell can correspond to
many plausible fine-scale temperature or precipitation patterns.
Conditional generative adversarial networks (GANs) were among the first
deep generative approaches used for climate and weather super-resolution.
They can sharpen fields and recover small-scale structure; examples
include adversarial super-resolution of wind and solar fields
\parencite{stengel2020adversarial_super_resolution}, precipitation
forecast correction and downscaling with CorrectorGAN
\parencite{price2022correctorgan}, and physically constrained GANs for
precipitation fields \parencite{hess2022physically_constrained_gan}. Their
drawbacks are familiar: adversarial training may be unstable,
probabilistic calibration is difficult, and physically meaningful
constraints must be added carefully.

Diffusion and score-based models provide a more explicit probabilistic
framework. Denoising diffusion models learn to reverse a noising process,
while score-based stochastic differential equation formulations learn a
time-dependent score field that transports noise into data
\parencite{ho2020ddpm,song2021score_sde}. In conditional downscaling, the
model samples \(y_H\) from noise while conditioning on \(x_L\), static
covariates, and possibly an initial deterministic estimate. The study
\emph{Generative Diffusion-based Downscaling for Climate} demonstrates
this idea in an idealized ERA5 experiment, recovering \(0.25^\circ\)
fields from \(2^\circ\) inputs and showing better fine-scale accuracy
than a standard U-Net, especially in spectral diagnostics
\parencite{watt2024diffusion_downscaling}. This is important because
spectral performance measures whether the model restores the spatial
scales that interpolation and deterministic regression tend to erase.

Recent diffusion papers also explore efficiency and physical realism.
Latent diffusion reduces computational cost by operating in a compressed
representation; \textcite{tomasi2025latent_diffusion_downscaling} use
this strategy to mimic kilometer-scale COSMO5.0\_CLM9 regional climate
simulations, asking whether AI can approximate aspects of dynamical
downscaling. Residual corrective diffusion models decompose the problem
into a deterministic mean prediction followed by a diffusion model for
the residual. This is the idea behind CorrDiff, which downscales
\(25\,\mathrm{km}\) reanalysis inputs to \(2\,\mathrm{km}\) regional
weather-model outputs over Taiwan, sharpening fronts and typhoons while
recovering realistic spectra and distributions
\parencite{mardani2025corrdiff}. A related direction replaces long
diffusion sampling chains with consistency models, which can generate
high-resolution samples in fewer steps. \textcite{hess2025consistency_downscaling}
show that such models can perform fast, scale-adaptive, uncertainty-aware
ESM downscaling and can generalize across Earth system model inputs
without retraining.

\subsubsection{Flow matching as an emerging alternative}

Flow matching is a recent generative paradigm that learns a
continuous-time velocity field transporting a simple base distribution
into the data distribution \parencite{lipman2023flow_matching}. For
downscaling, it has two attractive features. First, inference can be
formulated as solving an ordinary differential equation rather than
running many stochastic denoising steps, which may reduce sampling cost.
Second, physical constraints such as re-aggregation consistency, water
conservation, or humidity bounds can be expressed as penalties or
projection steps along the learned flow.

The first climate and weather applications are appearing quickly. A 2026
study uses residual flow matching to super-resolve coarse-resolution
weather forecasts, training only on reanalysis data and evaluating
whether re-coarsened outputs remain consistent with the original forecast
trajectory \parencite{delefosse2026flow_matching_weather}. Another 2026
study proposes physics-constrained adaptive flow matching for climate
downscaling, adding soft conservation constraints for precipitation and
humidity and evaluating a \(63\,\mathrm{km}\) to \(6.3\,\mathrm{km}\)
Central Europe task across multiple variables
\parencite{debeire2026physics_constrained_flow_matching}. These papers are
new, but they suggest that flow matching may be a useful successor or
complement to diffusion when fast sampling, ensemble generation, and
physical consistency must be balanced.

\subsubsection{Positioning of the present paper}
\label{sec:related-diagnostics}

The literature reviewed above is almost entirely \emph{corrective}: it
asks how to build a model that maps ERA5 (or another coarse field) onto
station-quality temperature. It is largely silent on the simpler, prior,
\emph{descriptive} question that motivates this paper: where and under
what conditions is the raw reanalysis-station gap largest, and which
easily available covariates predict it? Answering that question well is
what determines, in practice, which covariates a downscaling network such
as AMUN should be given, and where a residual-diffusion or
flow-matching correction such as CorrDiff is likely to matter most. To
our knowledge, no prior work provides a unified, Europe-wide,
covariate-level diagnostic of station--rescaled-product residuals with
uncertainty quantification that accounts for spatial clustering of
stations; we view the two literatures as complementary, with this paper
supplying the covariate-level diagnosis that the downscaling literature
currently takes largely as given.

% ===========================================================================
\section{Data, hackathon design, and methodology}
\label{sec:data}

\subsection{The GenHack'25 hackathon}

% TODO (Saad): confirm exact numbers (participants, teams, regions,
% dates) and add them here.
GenHack'25 was organised at \'Ecole Polytechnique, with Kayrros
providing technical support on data access and the geospatial processing
pipeline. Several hundred students, organised in teams, were given a
common task: identify and characterise the discrepancy between ERA5 2m
temperature and co-located weather-station observations. Teams were free
to choose their region of interest and their modelling approach; in
practice, most teams focused on a specific European country or
sub-region rather than on the continent as a whole. This paper does not
report on individual team submissions; instead, it distils the
covariates that recurred as explanatory across teams into a single,
consistent pipeline, and re-runs that pipeline over the whole of the
teams' combined domain, i.e.\ continental Europe.

\subsection{Data}

\begin{itemize}
\item \textbf{Target.} Station-observed 2m air temperature from
% TODO (Saad): name the station network(s) used, e.g. ECA\&D / national
% met services, and the time period covered.
[TK: station network], matched to the nearest ERA5 grid cell.
\item \textbf{Reference.} ERA5 reanalysis 2m temperature
\parencite{TODO:era5_hersbach2020} at the matched grid cell and
timestamp. [TK: confirm whether the ``rescaled'' product used is raw
ERA5, ERA5-Land, or a Kayrros-provided bias-adjusted product -- the
three drafts of this paper disagreed on this point and it should be
stated unambiguously once, here.]
\item \textbf{Residual.} $r_{it} = T^{\text{station}}_{it} -
T^{\text{ERA5}}_{it}$ for station $i$ at date $t$, so that $r_{it}>0$
means ERA5 (or the rescaled product) \emph{underestimates} station
temperature.
\item \textbf{Covariates.} Station temperature, precipitation (ERA5,
mm/day), wind speed (ERA5, m/s), NDVI, distance to the sea (km),
elevation (m), and DEGURBA degree-of-urbanisation class (water; very
low density rural; low density rural; rural cluster; suburban;
semi-dense urban cluster; dense urban cluster; urban centre).
\end{itemize}

\begin{table}[htbp]
\centering
\small
\caption{Descriptive statistics for continuous variables.}
\label{tab:desc_continuous}
\begin{tabular}{lrrrrrrrr}
\toprule
Variable & Nulls & Mean & Std & P05 & P25 & P50 & P75 & P95 \\
\midrule
era5\_celsius & 0 & 14 & 9.7 & -2.2 & 6.7 & 14 & 20 & 29 \\
residual & 0 & 1.6 & 2.5 & -2.1 & 0.25 & 1.4 & 2.7 & 5.8 \\
wind & 0 & 2.2 & 1.7 & 0.35 & 1.00 & 1.9 & 3.0 & 5.5 \\
dist\_to\_sea\_km & 0 & 83 & 94 & 0.21 & 8.4 & 47 & 130 & 280 \\
elevation & 0 & 340 & 420 & 4.0 & 42 & 180 & 500 & 1100 \\
ndvi & 0 & 0.39 & 0.25 & -0.024 & 0.23 & 0.41 & 0.58 & 0.75 \\
\bottomrule
\end{tabular}
\end{table}

\begin{table}[htbp]
\centering
\small
\caption{Distribution across DEGURBA urbanization classes.}
\label{tab:desc_degurba}
\begin{tabular}{llrrrr}
\toprule
Code & Label & Obs Count & Obs (\%) & Station Count & Station (\%) \\
\midrule
10 & Water & 152,911 & 1.6\% & 79 & 1.4\% \\
11 & Very low density rural & 3,514,982 & 36\% & 2,110 & 37\% \\
12 & Low density rural & 2,912,790 & 30\% & 1,622 & 29\% \\
13 & Rural cluster & 774,427 & 7.9\% & 450 & 7.9\% \\
21 & Suburban & 1,002,786 & 10.0\% & 550 & 9.7\% \\
22 & Semi-dense urban cluster & 156,261 & 1.6\% & 91 & 1.6\% \\
23 & Dense urban cluster & 512,965 & 5.2\% & 313 & 5.5\% \\
30 & Urban centre & 829,059 & 8.4\% & 449 & 7.9\% \\
\bottomrule
\end{tabular}
\end{table}

\begin{table}[htbp]
\centering
\small
\caption{Top 10 countries by weather station coverage.}
\label{tab:desc_country_top10}
\begin{tabular}{lrrrr}
\toprule
Country & Station Count & Station (\%) & Obs Count & Obs (\%) \\
\midrule
Italy & 1,047 & 18\% & 1,365,458 & 14\% \\
Spain & 923 & 16\% & 1,777,964 & 18\% \\
Norway & 703 & 12\% & 1,266,064 & 13\% \\
Germany & 655 & 12\% & 1,317,796 & 13\% \\
Austria & 597 & 11\% & 1,008,176 & 10.0\% \\
Sweden & 424 & 7.5\% & 836,831 & 8.5\% \\
Finland & 292 & 5.2\% & 573,328 & 5.8\% \\
Hungary & 166 & 2.9\% & 314,484 & 3.2\% \\
Poland & 145 & 2.6\% & 266,229 & 2.7\% \\
Ukraine & 115 & 2.0\% & 75,371 & 0.76\% \\
\bottomrule
\end{tabular}
\end{table}

\begin{table}[htbp]
\centering
\caption{Data summary: number of stations, observations, time span, and covariate coverage.}
\label{tab:data-summary}
\begin{tabular}{lc}
\toprule
Quantity & Value \\
\midrule
Number of stations & 5,664 \\
Number of station-days & 9,856,181 \\
Time period & 2020-01-01 to 2025-08-31 \\
Countries covered & 34 \\
Residual $r_{it}$, mean (sd) & 1.6 (2.5) $^\circ$C \\
\bottomrule
\end{tabular}
\end{table}

\subsection{Methodology}
\label{sec:methods}

\paragraph{Covariate profiles.} For each covariate $X_j$, we estimate the
partial dependence of the mean residual $\mathbb{E}[r \mid X_j]$ by
binning $X_j$ into quantiles and computing the mean residual within each
bin. This is a descriptive device: it does not impose a functional form
and is directly interpretable, at the cost of not controlling for other
covariates -- a limitation we return to in Section~\ref{sec:discussion}.
% TODO (Saad): confirm whether this is a *univariate* binned mean, or a
% partial-dependence profile extracted from a fitted multivariate model
% (e.g. gradient boosting) marginalising over the other covariates. The
% two have different causal interpretations and the methods section must
% say which one was used.

\paragraph{Inference.} Because residuals from the same station are
serially correlated and stations are spatially clustered, uncertainty is
quantified with a cluster bootstrap: stations (not observations) are
resampled with replacement, $B=1000$ replications, and we report 95\%
percentile confidence intervals for all profiled means and class means.
Clustering the resampling by station, rather than by observation,
respects the two-level structure of the data (dates nested in stations)
without imposing a parametric correlation model. The same procedure is
applied to the DEGURBA class breakdown (Section~\ref{sec:results-degurba}).

\paragraph{Interactions.} For the temperature$\times$wind interaction
(Section~\ref{sec:results-interaction}), observations are additionally
stratified into wind-speed terciles ([TK: confirm cut points, e.g.\ low
$\le 1.2\,\mathrm{m/s}$; mid $1.2$--$2.4\,\mathrm{m/s}$; high
$>2.4\,\mathrm{m/s}$]) before the temperature partial-dependence profile
is recomputed within each tercile.

% ===========================================================================
\section{Results}
\label{sec:results}

\subsection{ERA5 underestimates station temperature on average}
\label{sec:results-map}

\begin{figure}[htbp]
\centering
\includegraphics[width=0.85\textwidth]{european-map.png}
\caption{Mean residual per station, $\bar r_i = T_i^{-1}\sum_t r_{it}$,
all dates pooled. Positive (red) values indicate that the rescaled
product underestimates station temperature.}
\label{fig:map}
\end{figure}

Across the pooled European station network, the mean residual is
positive: station temperature exceeds the co-located ERA5 temperature on
average. This is not a local artefact: the bias is spatially widespread
rather than confined to a handful of stations or countries, with 90\% of stations showing a positive mean residual and an unweighted
sample mean of $\bar r_i \approx 1.5^\circ$C, though the magnitude
varies substantially by region -- largest in Albania, Bosnia and Herzegovina, Montenegro,
smallest, occasionally negative, in Russian Federation, Denmark, Netherlands.(Figure~\ref{fig:map}).

\subsection{Station temperature itself is a first-order explanatory
variable}
\label{sec:results-temperature}

\begin{figure}[htbp]
\centering
\includegraphics[width=\textwidth]{sensitivities.png}
\caption{Partial dependence of the mean residual on each of six
covariates: station temperature, precipitation, wind speed, NDVI,
distance to the sea, and elevation. Covariates are binned into
quantiles; shaded bands are 95\% confidence intervals from a
station-clustered bootstrap ($B=1000$).}
\label{fig:pdp}
\end{figure}

The partial dependence of the mean residual on station temperature is
close to monotonically increasing (Figure~\ref{fig:pdp}, top left): the
underestimation is smallest at low station temperatures and largest at
high station temperatures, rising from approximately $0.84\,^\circ$C at
$-20\,^\circ$C to $2.3\,^\circ$C at $+40\,^\circ$C. In other words, the
rescaled product's cold bias worsens precisely on the warmest days --
the opposite of what one would want for downstream applications
sensitive to heat extremes (heatwave detection, energy demand, crop
stress).

\subsection{Precipitation and NDVI show at most a minor effect}
\label{sec:results-precip}

The partial dependence of the residual on precipitation is close to flat, with confidence
bands that overlap substantially across the observed range (0.001
–1.3 mm/day, with mean residual hovering between 1.5 and 1.6~\text{°C}, 95\% CI [1.5, 1.7]~\text{°C}) (Figure 2, top centre). We therefore report no more than a minor precipitation effect at
continental scale – a statement about the observed range rather than an absence of
physical effect. The partial dependence on NDVI is likewise close to flat across its observed range (-0.40–0.78, mean residual between 1.5 and 1.6~\text{°C}), with substantially overlapping confidence bands (Figure 2, bottom left);
any effect is below the detection threshold of the current sample.
% TODO (Saad): decide between "no detectable effect" and "small effect"
% based on whether a formal test (e.g. comparing bin means against a
% null of no association) supports a genuine, if modest, relationship;
% see Section~\ref{sec:discussion}. Do not let this wording drift
% between sections in the final draft.

\subsection{Wind speed strongly damps the underestimation}
\label{sec:results-wind}

Wind speed is the covariate with the clearest monotonic marginal
association in the panel (Figure~\ref{fig:pdp}, top right): the residual
falls steadily as wind speed increases, from its largest values at
near-calm conditions to roughly half that magnitude at the highest
observed wind speeds (from $\sim 2.1^\circ$C to $\sim 1.1^\circ$C).

\subsection{Distance to the sea and elevation}
\label{sec:results-geo}

The residual decreases with distance to the sea beyond a short coastal
band, though the profile is not monotonic over its full range: mean
underestimation initially rises from 1.69°C (at ~2 km) to 1.92°C (at ~17 km)
before declining steadily further inland to $1.15^\circ$C at $\sim 364$ km
(Figure~\ref{fig:pdp}, bottom centre).
Elevation shows the expected sign given prior regional findings
\parencite{satyapragyan2026bias,sebbar2023ml_downscaling}: the
underestimation increases with elevation (rising from $1.34^\circ$C at $\sim 8$ m to $1.65^\circ$C at
$\sim 352$ m), consistent with the rescaled
product's coarse grid cells averaging over sub-grid topographic
variability that individual mountain stations do not experience
(Figure~\ref{fig:pdp}, bottom right).
% TODO (Saad): flag and, if possible, explain any near-coast feature
% rather than smoothing over it -- check station counts per bin, since
% it may be a sample-size artefact rather than a genuine coastal effect.


\subsection{Degree of urbanisation}
\label{sec:results-degurba}

\begin{figure}[htbp]
\centering
\includegraphics[width=0.85\textwidth]{images/degurba.png}
\caption{Mean residual per DEGURBA urbanisation class (top), with
station share per class (bottom). Error bars are 95\% station-clustered
bootstrap confidence intervals.}
\label{fig:degurba}
\end{figure}

Stations classified as semi-dense or dense urban clusters show
noticeably higher mean residuals than stations classified as water or
very-low-density rural, with rural-cluster and suburban stations
occupying an intermediate position; the residual peaks in semi-dense
urban clusters at approximately $+2.55^\circ$C and remains elevated in
the densest classes (Figure~\ref{fig:degurba}). Because the residual is
positive in \emph{all} classes, urbanisation amplifies but does not
solely explain the underestimation. Two readings are compatible with
this pattern: the rescaled product under-represents the urban heat
signal measured by stations, and/or station siting in built environments
warms the reference relative to the grid-cell average. Our design cannot
separate these two components; we flag this as a priority for follow-up
work. Together with the elevation result, this points to two largely
distinct geographic drivers of the residual: topographic (elevation,
distance to sea) and land-use / built-environment (urbanisation class).

\subsection{A temperature$\times$wind interaction}
\label{sec:results-interaction}

\begin{figure}[htbp]
\centering
\includegraphics[width=0.8\textwidth]{joined-effects.png}
\caption{Partial dependence of the mean residual on station temperature,
stratified by wind-speed tercile. Shaded bands are 95\% station-clustered
bootstrap confidence intervals.}
\label{fig:interaction}
\end{figure}

Stratifying the station-temperature partial dependence profile by
wind-speed tercile reveals that the temperature effect documented in
Section~\ref{sec:results-temperature} is not uniform
(Figure~\ref{fig:interaction}): under low wind, the underestimation
rises sharply with station temperature and reaches its largest values in
the sample; under high wind, the same temperature range produces a much
flatter, substantially smaller residual. At low station temperatures the
three wind terciles are close to indistinguishable; the gap opens up
specifically as temperature rises, reaching $\sim 1.8^\circ$C
under calm conditions versus $\sim 2.8^\circ$C under strong wind] at the
warmest temperatures. This is consistent with a physical mechanism in
which low wind favours a stable near-surface boundary layer that the
rescaled product's coarser vertical/horizontal resolution does not fully
resolve, so that the model's warm-day bias is largest precisely when
vertical mixing is weakest. The amplification of the bias by stagnant
boundary layers is, jointly with the urbanisation gradient, the most
physical signal in the data: both point to unresolved surface-layer
processes rather than to a uniform calibration offset.


\begin{table}[htbp]
\centering
\caption{Summary of marginal associations between covariates and the residual.}
\label{tab:results-summary}
\resizebox{\linewidth}{!}{%
\begin{tabular}{lccc}
\toprule
Covariate & Direction of association & Approx.\ range ($^\circ$C) & Detectable? \\
\midrule
Station temperature & Increasing & 0.84--2.3 & Yes (strong) \\
Wind speed & Decreasing & 1.1--2.1 & Yes (strong) \\
Distance to sea & Decreasing (non-monotonic near coast) & 1.1--1.9 & Yes \\
DEGURBA (urban $>$ rural) & Increasing with density & 1.1--2.55 & Yes (weak)\\
Elevation & Increasing & 1.3--1.7 & Yes (weaker) \\
Precipitation & None / negligible & 1.5--1.6 & No \\
NDVI & None / negligible & 1.5--1.6 & No \\
\bottomrule
\end{tabular}}
\end{table}

% ===========================================================================
\section{Discussion}
\label{sec:discussion}

\paragraph{Interpretation and threats to validity.}
Several caveats qualify the results above. \emph{(i) Magnitude.} A mean
underestimation on the order of [TK]$^\circ$C is large relative to
published reanalysis biases over Europe; part of it may reflect
processing artefacts rather than genuine product error -- elevation
mismatch between station and grid cell, differing averaging windows, or
station screening. We regard artefact checks as a precondition for
downstream use of these estimates [TK: report artefact checks, or move
to an appendix]. \emph{(ii) Univariate profiles and confounding by
pooling.} The results in Section~\ref{sec:results} are, as currently
implemented, marginal (univariate) associations between the residual and
each covariate taken separately. Several of these covariates are
themselves correlated -- mountain stations tend to be further from the
sea and less urbanised, for instance -- so a univariate profile can
partly reflect the influence of a correlated covariate rather than the
one being plotted. A joint model (e.g.\ a multivariate regression, or a
fitted gradient-boosting model with true multivariate partial
dependence) would let each covariate's effect be interpreted holding the
others fixed, which is a meaningfully stronger claim than what
Table~\ref{tab:results-summary} currently supports. All profiles also
pool dates, so the temperature profile in particular partly reflects the
seasonal cycle; station fixed effects or deseasonalised residuals would
isolate within-station variation [TK: add a mixed-model specification if
time allows]. \emph{(iii) Reference ambiguity.} Station thermometers and
grid cells average different spatial footprints; over cities and coasts,
part of the residual measures representativeness rather than accuracy.
\emph{(iv) Statistical significance of ``minor'' effects.} The
precipitation and NDVI effects are visually small and their confidence
bands largely overlap across bins; before describing either as
``minor but noticeable'' in the final text, a formal test against a null
of no association is advisable.

More substantively, the two strongest results in this analysis -- the
wind-speed effect and the temperature$\times$wind interaction -- deserve
to be stated as clearly among the paper's central conclusions
(Section~\ref{sec:intro}), since they are, on the evidence gathered here,
at least as strong as the elevation effect that originally motivated
this line of work. The interaction result in particular offers a
candidate physical mechanism (near-surface stability under low wind)
that ties several of the other findings together and gives the paper a
sharper, more specific contribution than a list of marginal associations.

\paragraph{Consequences for downscaling.}
The diagnostics of Section~\ref{sec:results} define a natural empirical
target for downscaling research in the sense of
Section~\ref{sec:related-diagnostics}: the conditional structure of the
residual -- its dependence on temperature, wind, elevation, distance to
sea, and urban class, and its spatial clustering -- is precisely the
component that a residual-based or bias-aware downscaling model should
learn. The strong temperature$\times$wind interaction further suggests
that dynamic covariates (boundary-layer state) carry information that
static covariate sets (topography, land cover, NDVI) do not. Finally,
our results caution against treating rescaled products as ground truth
in the training and evaluation of downscaling models: doing so transfers
the biases documented here into the learned model.

% ===========================================================================
\section{Conclusion}
\label{sec:conclusion}

Using the collective output of the GenHack'25 hackathon recoded into a
single European pipeline, we characterised the errors of a rescaled
reanalysis temperature product against weather stations. The product
underestimates station temperature on average; the underestimation grows
with temperature itself, with elevation, and towards the coast, and is
strongly amplified by low wind and by urbanisation. Precipitation and
vegetation show at most a minor effect at continental scale. These
conclusions delineate where rescaled temperatures can be trusted and
provide an empirical error model for downstream uncertainty
quantification. Users of ERA5-type temperature in downstream models
should, at minimum, be aware that the error grows with temperature and
shrinks with wind speed; future downscaling or bias-correction models
(Section~\ref{sec:related-ml}) should treat wind speed and DEGURBA class
as first-class conditioning covariates alongside the more commonly used
elevation. Planned extensions include deseasonalised within-station
inference, artefact controls, a joint (multivariate) specification of
the covariate effects, and the separation of station-siting from product
bias in urban areas.

\printbibliography

\end{document}
